\section{Chapter 8}

\begin{verse}
The transcendent quality resembles water\\
Without resisting it assumes the form of each thing\\
It takes the lowest position that men scorn (because they are seeking the heights)\\
The farther one is from acting in common, the more one is close to the Way\\
Thus the True Man in the material domain holds remaining in the place he has as good\\
In the domain of feeling he holds the depth of the abyss as good\\
If he gives, he holds (impersonal) generosity as good\\
Speaking, he holds the truth as good\\
Governing he holds ordered development as good\\
Acting, he holds the watchful realization of the end as good\\
In the practical domain, he holds being at the right place as good\\
Really, even thanks to adapting without conflict\\
Nothing changes his being
\end{verse}

Classical Taoist comparison of water for subtle acting and the ``higher virtue'' (\emph{shang shan}). Adapting himself exteriorly to things and situations, he does not exalt but descends below (which means: beyond superficiality) in order to exercise the desired influence. According to another aspect: in this context the saying ``fracture but do not bend'' is tied to the mania of the I; ``bending myself (=water) it fractures'' could be contrasted. The whole of the behavior illustrated in this chapter is designated as \emph{I hsing} (which is equivalent to ``ease of nature''). The positive counterpart is also recommended: a spontaneous and essential mode of working, beyond opposites.

\paragraph{Chinese text and literal translation}


\textbf{Chapter 8 (\begin{CJK*}{UTF8}{bsmi}第八章\end{CJK*})}

\columnratio{.4}
\begin{paracol}{2}
\begin{verse}
\begin{CJK*}{UTF8}{bsmi}
上善若水。\\
水善利萬物而有靜，\\
處眾人之所惡，\\
故幾於道。\\
居善地，\\
心善淵，\\
與善仁，\\
言善信，\\
正善治，\\
事善能，\\
動善時。\\
夫唯不爭，\\
故無尤。
\end{CJK*}
\end{verse}
\switchcolumn
\begin{verse}
The best quality/character is like water.\\
The water's goodness is that it benefits the myriad things but does not quarrel,\\
and it willingly goes to where others hate,\\
Thus it is almost like the Tao.\\
It is good to be/live on the ground,\\
to deepen a heart,\\
to love people while associating with them.\\
to keep one's word while talking,\\
to be peace while governing,\\
to do what one is capable of,\\
to act at a fit time.\\
Because of the non-fighting-over,\\
there will be no blame.
\end{verse}
\end{paracol}


\flrightit{Posted on 2020-09-20 by Lao Tzu}

